\import{algebra/module/free.tex}
\import{algebra/module/span.tex}
\import{algebra/module/linear-map.tex}

\section{Indexed spanning families and bases}

An indexed family determines a linear map out of the free module by sending
each canonical basis vector to the corresponding member of the family.
Spanning, linear independence, and being a basis say respectively that this
map is surjective, injective, and bijective.  The structural and typing guards
are part of each predicate, so these phrases remain meaningful as global
notions.  The permitted trivial ring is genuinely degenerate here: every
module carrier and every free-module carrier is a singleton, so every
well-typed indexed family is a basis, for every index set.

\begin{definition}\label{linear_combination_map}
    $\linearCombinationMap{R}{I}{M}{v}
        =\freeModuleLift{R}{I}{M}{v}$.
\end{definition}

\begin{definition}\label{indexed_module_spanning}
    $v$ spans $M$ over $R$ indexed by $I$ iff
        $R$ is a ring and $M$ is a left module over $R$ and
        $v\in\funs{I}{\carrier[M]}$ and
        $\linearCombinationMap{R}{I}{M}{v}$ is a module surjection from
            $\FreeModule{R}{I}$ to $M$ over $R$.
\end{definition}

\begin{definition}\label{indexed_module_linear_independence}
    $v$ is linearly independent in $M$ over $R$ indexed by $I$ iff
        $R$ is a ring and $M$ is a left module over $R$ and
        $v\in\funs{I}{\carrier[M]}$ and
        $\linearCombinationMap{R}{I}{M}{v}$ is injective.
\end{definition}

\begin{definition}\label{indexed_module_basis}
    $v$ forms a basis of $M$ over $R$ indexed by $I$ iff
        $R$ is a ring and $M$ is a left module over $R$ and
        $v\in\funs{I}{\carrier[M]}$ and
        $\linearCombinationMap{R}{I}{M}{v}$ is a module isomorphism from
            $\FreeModule{R}{I}$ to $M$ over $R$.
\end{definition}

\begin{proposition}\label{linear_combination_map_linear}
    Let $R$ be a ring.
    Let $M$ be a left module over $R$.
    Let $v\in\funs{I}{\carrier[M]}$.
    Then $\linearCombinationMap{R}{I}{M}{v}$ is a linear map from
        $\FreeModule{R}{I}$ to $M$ over $R$.
\end{proposition}
\begin{proof}
    Follows by \auto{linear_combination_map,free_module_lift_linear}.
\end{proof}

\begin{proposition}\label{linear_combination_map_basis_apply}
    Let $R$ be a ring.
    Let $M$ be a left module over $R$.
    Let $v\in\funs{I}{\carrier[M]}$ and $i\in I$.
    Then
        $\apply{\linearCombinationMap{R}{I}{M}{v}}{
            \freeBasisVector{R}{I}{i}}=v(i)$.
\end{proposition}
\begin{proof}
    Follows by
        \auto{linear_combination_map,free_module_lift_basis_apply}.
\end{proof}

\begin{proposition}\label{indexed_module_spanning_guards}
    Suppose $v$ spans $M$ over $R$ indexed by $I$.
    Then $R$ is a ring and $M$ is a left module over $R$ and
        $v\in\funs{I}{\carrier[M]}$.
\end{proposition}
\begin{proof}
    Follows by \obv{indexed_module_spanning}.
\end{proof}

\begin{proposition}\label{indexed_module_linear_independence_guards}
    Suppose $v$ is linearly independent in $M$ over $R$ indexed by $I$.
    Then $R$ is a ring and $M$ is a left module over $R$ and
        $v\in\funs{I}{\carrier[M]}$.
\end{proposition}
\begin{proof}
    Follows by \obv{indexed_module_linear_independence}.
\end{proof}

\begin{proposition}\label{indexed_module_basis_guards}
    Suppose $v$ forms a basis of $M$ over $R$ indexed by $I$.
    Then $R$ is a ring and $M$ is a left module over $R$ and
        $v\in\funs{I}{\carrier[M]}$.
\end{proposition}
\begin{proof}
    Follows by \obv{indexed_module_basis}.
\end{proof}

\begin{theorem}\label{indexed_module_basis_iff}
    Let $R$ be a ring.
    Let $M$ be a left module over $R$.
    Let $v\in\funs{I}{\carrier[M]}$.
    Then $v$ forms a basis of $M$ over $R$ indexed by $I$ iff
        $v$ spans $M$ over $R$ indexed by $I$ and
        $v$ is linearly independent in $M$ over $R$ indexed by $I$.
\end{theorem}
\begin{proof}
    Let $F=\FreeModule{R}{I}$ and
        $f=\linearCombinationMap{R}{I}{M}{v}$.
    $f$ is a linear map from $F$ to $M$ over $R$ by
        \auto{linear_combination_map_linear}.
    Show that if $v$ forms a basis of $M$ over $R$ indexed by $I$ then
        $f$ is a module isomorphism from $F$ to $M$ over $R$.
    \begin{subproof}
        Assume $v$ forms a basis of $M$ over $R$ indexed by $I$.
        $\linearCombinationMap{R}{I}{M}{v}$ is a module isomorphism from
            $\FreeModule{R}{I}$ to $M$ over $R$ by
            \obv{indexed_module_basis}.
    \end{subproof}
    If $f$ is a module isomorphism from $F$ to $M$ over $R$ then
        $f$ is a module surjection from $F$ to $M$ over $R$ and
        $f$ is injective by
        \auto{module_isomorphism_is_surjection,
            module_isomorphism_is_bijection,bijections}.
    If $f$ is a module surjection from $F$ to $M$ over $R$ and
        $f$ is injective then
        $f\in\bijections{\carrier[F]}{\carrier[M]}$ by
        \obv{module_surjection_is_surj,bijections}.
    Hence if $f$ is a module surjection from $F$ to $M$ over $R$ and
        $f$ is injective then
        $f$ is a module isomorphism from $F$ to $M$ over $R$ by
        \obv{module_isomorphism}.
    Follows by
        \auto{indexed_module_spanning,
            indexed_module_linear_independence,indexed_module_basis}.
\end{proof}

\begin{theorem}\label{trivial_ring_every_indexed_family_is_basis}
    Let $R$ be a ring.
    Let $M$ be a left module over $R$.
    Suppose $\ringZero[R]=\ringOne[R]$.
    Let $I$ be a set and $v\in\funs{I}{\carrier[M]}$.
    Then $v$ forms a basis of $M$ over $R$ indexed by $I$.
\end{theorem}
\begin{proof}
    Let $F=\FreeModule{R}{I}$ and
        $f=\linearCombinationMap{R}{I}{M}{v}$.
    $F$ is an exact left module over $R$ by
        \auto{free_module_is_exact}.
    Hence $F$ is a left module over $R$.
    $\carrier[F]=\{\neutral[F]\}$ and
        $\carrier[M]=\{\neutral[M]\}$ by
        \obv{trivial_ring_module_carrier_singleton}.
    $f$ is a linear map from $F$ to $M$ over $R$ by
        \auto{linear_combination_map_linear}.
    Hence $f\in\funs{\carrier[F]}{\carrier[M]}$ by
        \obv{linear_map_type}.
    Hence $f$ is a function and $\dom{f}=\carrier[F]$ by
        \obv{funs_is_function,funs}.
    Show that $f\in\surj{\carrier[F]}{\carrier[M]}$.
    \begin{subproof}
        Show that for every $y\in\carrier[M]$ there exists
            $x\in\carrier[F]$ such that $f(x)=y$.
        \begin{subproof}
            Fix $y\in\carrier[M]$.
            $y=\neutral[M]$ by \obv{singleton_iff}.
            $\neutral[F]\in\carrier[F]$ by
                \obv{module_is_group,group_is_monoid,
                    monoid_is_unital_magma,unitalmagma_type}.
            $f(\neutral[F])=\neutral[M]$ by \obv{linear_map_zero}.
        \end{subproof}
        Follows by \obv{surj}.
    \end{subproof}
    Show that $f$ is injective.
    \begin{subproof}
        Show that for all $x,y\in\dom{f}$ such that $f(x)=f(y)$
            we have $x=y$.
        \begin{subproof}
            Fix $x,y\in\dom{f}$.
            Assume $f(x)=f(y)$.
            $x,y\in\{\neutral[F]\}$.
            Follows by \auto{singleton_iff}.
        \end{subproof}
        Follows by \obv{injective_function}.
    \end{subproof}
    Hence $f\in\bijections{\carrier[F]}{\carrier[M]}$ by
        \obv{bijections}.
    Hence $f$ is a module isomorphism from $F$ to $M$ over $R$ by
        \obv{module_isomorphism}.
    Follows by \obv{indexed_module_basis}.
\end{proof}

\subsection{Span and the linear-combination map}

\begin{theorem}\label{indexed_family_span_eq_linear_combination_range}
    Let $R$ be a ring.
    Let $M$ be a left module over $R$.
    Let $v\in\funs{I}{\carrier[M]}$.
    Then
        $\moduleSpan{R}{M}{\ran{v}}
            =\ran{\linearCombinationMap{R}{I}{M}{v}}$.
\end{theorem}
\begin{proof}
    Let $S=\FreeModule{R}{I}$ and
        $f=\linearCombinationMap{R}{I}{M}{v}$.
    Let $P=\moduleSpan{R}{M}{\ran{v}}$ and
        $Q=\InducedSubmodule{R}{M}{P}$.
    $\ran{v}\subseteq\carrier[M]$ by \obv{funs_ran}.
    Hence $P$ is a submodule of $M$ over $R$ by
        \auto{module_span_is_submodule}.
    $f$ is a linear map from $S$ to $M$ over $R$ by
        \auto{linear_combination_map_linear}.
    Hence $\moduleImage{f}{S}$ is a submodule of $M$ over $R$ by
        \obv{module_image_is_submodule}.
    $f\in\funs{\carrier[S]}{\carrier[M]}$ by
        \obv{linear_map_type}.
    Hence $f$ is a function and $\dom{f}=\carrier[S]$ by
        \obv{funs_is_function,funs}.
    Hence $\ran{f}=\img{f}{\carrier[S]}$ by \auto{img_dom}.
    Hence $\ran{f}=\moduleImage{f}{S}$ by
        \obv{module_image}.
    Hence $\ran{f}$ is a submodule of $M$ over $R$.
    Show that $\ran{v}\subseteq\ran{f}$.
    \begin{subproof}
        Show that for every $y\in\ran{v}$ we have $y\in\ran{f}$.
        \begin{subproof}
            Fix $y\in\ran{v}$.
            $v$ is a function by \obv{funs_is_function}.
            Take $i\in\dom{v}$ such that $v(i)=y$ by
                \obv{fun_ran_iff}.
            Hence $i\in I$ by \obv{funs}.
            Let $b=\freeBasisVector{R}{I}{i}$.
            $b\in\carrier[S]$ by \auto{free_basis_vector_type}.
            Hence $b\in\dom{f}$.
            $f(b)=v(i)$ by
                \auto{linear_combination_map_basis_apply}.
            Hence $f(b)=y$.
            Follows by \auto{fun_ran_iff}.
        \end{subproof}
        Follows by \obv{subseteq}.
    \end{subproof}
    Hence $P\subseteq\ran{f}$ by \auto{module_span_least}.
    $Q$ is an exact left module over $R$ by
        \auto{induced_submodule_is_exact}.
    Hence $Q$ is a left module over $R$.
    $\carrier[Q]=P$ by \auto{induced_submodule_carrier}.
    $\ran{v}\subseteq P$ by
        \obv{module_span_contains_generators}.
    $v$ is a relation and $\dom{v}=I$ by
        \obv{funs_is_relation,funs}.
    Hence $\dom{v}\subseteq I$ by \obv{subseteq_refl}.
    $v\in\rels{I}{P}$ by
        \obv{rels_intro_dom_and_ran}.
    Hence $v\in\funs{I}{P}$ by \obv{funs}.
    Hence $v\in\funs{I}{\carrier[Q]}$.
    Let $k=\freeModuleLift{R}{I}{Q}{v}$.
    $k$ is a linear map from $S$ to $Q$ over $R$ by
        \auto{free_module_lift_linear}.
    $\moduleSubstructure{R}{Q}{M}$ by
        \auto{induced_submodule_is_module_substructure}.
    Hence $k$ is a linear map from $S$ to $M$ over $R$ by
        \obv{linear_map_codomain_enlargement}.
    For every $i\in I$ we have
        $k(\freeBasisVector{R}{I}{i})=v(i)$ by
        \auto{free_module_lift_basis_apply}.
    Hence $k=\freeModuleLift{R}{I}{M}{v}$ by
        \auto{free_module_lift_unique}.
    Hence $k=f$ by \obv{linear_combination_map}.
    $k\in\funs{\carrier[S]}{\carrier[Q]}$ by
        \auto{free_module_lift_type}.
    Hence $\ran{k}\subseteq\carrier[Q]$ by \obv{funs_ran}.
    Hence $\ran{f}\subseteq P$.
    Follows by \obv{subseteq_antisymmetric}.
\end{proof}

\subsection{Reindexing}

\begin{theorem}\label{linear_combination_map_reindex}
    Let $R$ be a ring.
    Let $M$ be a left module over $R$.
    Let $e\in\bijections{J}{I}$ and
        $v\in\funs{I}{\carrier[M]}$.
    Let $F=I\times\{\LeftRegularModule{R}\}$ and
        $A=\moduleAdditiveGroupFamily{I}{F}$.
    Let $r=\directSumGroupReindex{I}{J}{A}{e}$.
    Then
        $\linearCombinationMap{R}{J}{M}{v\circ e}\circ r
            =\linearCombinationMap{R}{I}{M}{v}$.
\end{theorem}
\begin{proof}
    Let $S=\FreeModule{R}{I}$ and $T=\FreeModule{R}{J}$.
    Let $g=\linearCombinationMap{R}{J}{M}{v\circ e}$ and
        $f=\linearCombinationMap{R}{I}{M}{v}$.
    $e\in\funs{J}{I}$ by \obv{bijections_to_funs}.
    Hence $v\circ e\in\funs{J}{\carrier[M]}$ by
        \obv{funs_circ}.
    $r$ is a module isomorphism from $S$ to $T$ over $R$ by
        \auto{bijective_free_module_reindex_isomorphism}.
    Hence $r$ is a linear map from $S$ to $T$ over $R$ by
        \obv{module_isomorphism_is_linear_map}.
    $g$ is a linear map from $T$ to $M$ over $R$ by
        \auto{linear_combination_map_linear}.
    Hence $g\circ r$ is a linear map from $S$ to $M$ over $R$ by
        \obv{linear_map_composition}.
    $r\in\funs{\carrier[S]}{\carrier[T]}$ and
        $g\in\funs{\carrier[T]}{\carrier[M]}$ by
        \obv{linear_map_type}.
    Hence $r$ and $g$ are functions by \auto{funs_is_function}.
    $g$ is composable with $r$ by \auto{funs_ran,funs}.
    Show that for every $i\in I$ we have
        $(g\circ r)(\freeBasisVector{R}{I}{i})=v(i)$.
    \begin{subproof}
        Fix $i\in I$.
        Let $b=\freeBasisVector{R}{I}{i}$.
        $e\in\surj{J}{I}$ by \obv{bijections}.
        Take $j\in J$ such that $e(j)=i$ by \obv{surj}.
        $r(\freeBasisVector{R}{I}{e(j)})
            =\freeBasisVector{R}{J}{j}$ by
            \obv{free_basis_vector_reindex}.
        Hence $r(b)
            =\freeBasisVector{R}{J}{j}$.
        $b\in\carrier[S]$ by
            \auto{free_basis_vector_type}.
        Hence $b\in\dom{r}$ by \auto{funs}.
        Hence $(g\circ r)(b)=g(r(b))$ by \obv{circ_apply}.
        $g(\freeBasisVector{R}{J}{j})=(v\circ e)(j)$ by
            \auto{linear_combination_map_basis_apply}.
        $v$ and $e$ are functions and $j\in\dom{e}$ by
            \auto{funs_is_function,funs}.
        $v$ is composable with $e$ by \auto{funs_ran,funs}.
        Hence $(v\circ e)(j)=v(e(j))$ by \obv{circ_apply}.
    \end{subproof}
    $g\circ r=\freeModuleLift{R}{I}{M}{v}$ by
        \auto{free_module_lift_unique}.
    Follows by \auto{linear_combination_map}.
\end{proof}

\begin{theorem}\label{indexed_module_spanning_reindex}
    Let $R$ be a ring.
    Let $M$ be a left module over $R$.
    Let $e\in\bijections{J}{I}$ and
        $v\in\funs{I}{\carrier[M]}$.
    Then $v$ spans $M$ over $R$ indexed by $I$ iff
        $v\circ e$ spans $M$ over $R$ indexed by $J$.
\end{theorem}
\begin{proof}
    Let $S=\FreeModule{R}{I}$ and $T=\FreeModule{R}{J}$.
    Let $F=I\times\{\LeftRegularModule{R}\}$ and
        $A=\moduleAdditiveGroupFamily{I}{F}$.
    Let $r=\directSumGroupReindex{I}{J}{A}{e}$.
    Let $f=\linearCombinationMap{R}{I}{M}{v}$ and
        $g=\linearCombinationMap{R}{J}{M}{v\circ e}$.
    $e\in\funs{J}{I}$ by \obv{bijections_to_funs}.
    Hence $v\circ e\in\funs{J}{\carrier[M]}$ by
        \obv{funs_circ}.
    $f$ is a linear map from $S$ to $M$ over $R$ and
        $g$ is a linear map from $T$ to $M$ over $R$ by
        \auto{linear_combination_map_linear}.
    Hence $f\in\funs{\carrier[S]}{\carrier[M]}$ and
        $g\in\funs{\carrier[T]}{\carrier[M]}$ by
        \obv{linear_map_type}.
    $r$ is a module isomorphism from $S$ to $T$ over $R$ by
        \auto{bijective_free_module_reindex_isomorphism}.
    Hence $r\in\bijections{\carrier[S]}{\carrier[T]}$ by
        \obv{module_isomorphism_is_bijection}.
    Hence $r\in\surj{\carrier[S]}{\carrier[T]}$ by
        \obv{bijections}.
    Hence $\ran{r}=\carrier[T]$ by \obv{ran_of_surj}.
    $\dom{g}=\carrier[T]$ by \obv{funs}.
    Hence $\dom{g}=\ran{r}$.
    $r$ is a function by \obv{bijection_is_function}.
    $g$ is a function by \auto{funs_is_function}.
    Hence $\ran{g\circ r}=\ran{g}$ by \obv{ran_circ_exact}.
    $g\circ r=f$ by \auto{linear_combination_map_reindex}.
    Hence $\ran{f}=\ran{g}$.
    Hence $f\in\surj{\carrier[S]}{\carrier[M]}$ iff
        $g\in\surj{\carrier[T]}{\carrier[M]}$ by
        \auto{funs_surj_iff}.
    Hence $f$ is a module surjection from $S$ to $M$ over $R$ iff
        $g$ is a module surjection from $T$ to $M$ over $R$ by
        \auto{module_surjection}.
    Follows by \auto{indexed_module_spanning}.
\end{proof}

\begin{theorem}\label{indexed_module_linear_independence_reindex}
    Let $R$ be a ring.
    Let $M$ be a left module over $R$.
    Let $e\in\bijections{J}{I}$ and
        $v\in\funs{I}{\carrier[M]}$.
    Then $v$ is linearly independent in $M$ over $R$ indexed by $I$ iff
        $v\circ e$ is linearly independent in $M$ over $R$ indexed by $J$.
\end{theorem}
\begin{proof}
    Let $S=\FreeModule{R}{I}$ and $T=\FreeModule{R}{J}$.
    Let $F=I\times\{\LeftRegularModule{R}\}$ and
        $A=\moduleAdditiveGroupFamily{I}{F}$.
    Let $r=\directSumGroupReindex{I}{J}{A}{e}$.
    Let $f=\linearCombinationMap{R}{I}{M}{v}$ and
        $g=\linearCombinationMap{R}{J}{M}{v\circ e}$.
    $e\in\funs{J}{I}$ by \obv{bijections_to_funs}.
    Hence $v\circ e\in\funs{J}{\carrier[M]}$ by
        \obv{funs_circ}.
    $f$ is a linear map from $S$ to $M$ over $R$ and
        $g$ is a linear map from $T$ to $M$ over $R$ by
        \auto{linear_combination_map_linear}.
    Hence $f\in\funs{\carrier[S]}{\carrier[M]}$ and
        $g\in\funs{\carrier[T]}{\carrier[M]}$ by
        \obv{linear_map_type}.
    Hence $f$ and $g$ are functions and
        $\dom{f}=\carrier[S]$ and $\dom{g}=\carrier[T]$ by
        \auto{funs_is_function,funs}.
    $r$ is a module isomorphism from $S$ to $T$ over $R$ by
        \auto{bijective_free_module_reindex_isomorphism}.
    Hence $r\in\bijections{\carrier[S]}{\carrier[T]}$ by
        \obv{module_isomorphism_is_bijection}.
    Hence $r$ is a function and $r$ is injective and
        $r\in\surj{\carrier[S]}{\carrier[T]}$ by
        \auto{bijections,bijection_is_function}.
    Hence $\dom{r}=\carrier[S]$ and $\ran{r}=\carrier[T]$ by
        \obv{bijections_dom,ran_of_surj}.
    Hence $\ran{r}=\dom{g}$.
    Hence $g$ is composable with $r$ by \obv{subseteq_refl}.
    $g\circ r=f$ by \auto{linear_combination_map_reindex}.
    Show that if $g$ is injective then $f$ is injective.
    \begin{subproof}
        Assume $g$ is injective.
        Hence $g\circ r$ is injective by \obv{circ_injective}.
    \end{subproof}
    Show that if $f$ is injective then $g$ is injective.
    \begin{subproof}
        Assume $f$ is injective.
        Show that for all $x,y\in\dom{g}$ such that $g(x)=g(y)$
            we have $x=y$.
        \begin{subproof}
            Fix $x,y\in\dom{g}$.
            Assume $g(x)=g(y)$.
            Hence $x,y\in\carrier[T]$.
            Take $a\in\carrier[S]$ such that $r(a)=x$ by
                \obv{surj}.
            Take $b\in\carrier[S]$ such that $r(b)=y$ by
                \obv{surj}.
            Hence $a,b\in\dom{r}$.
            Hence $(g\circ r)(a)=g(r(a))$ and
                $(g\circ r)(b)=g(r(b))$ by \obv{circ_apply}.
            Hence $f(a)=f(b)$.
            $a,b\in\dom{f}$.
            Hence $a=b$ by \auto{injective_function}.
        \end{subproof}
        Follows by \auto{injective_function}.
    \end{subproof}
    Hence $f$ is injective iff $g$ is injective by \auto.
    Follows by \auto{indexed_module_linear_independence}.
\end{proof}

\begin{theorem}\label{indexed_module_basis_reindex}
    Let $R$ be a ring.
    Let $M$ be a left module over $R$.
    Let $e\in\bijections{J}{I}$ and
        $v\in\funs{I}{\carrier[M]}$.
    Then $v$ forms a basis of $M$ over $R$ indexed by $I$ iff
        $v\circ e$ forms a basis of $M$ over $R$ indexed by $J$.
\end{theorem}
\begin{proof}
    $e\in\funs{J}{I}$ by \obv{bijections_to_funs}.
    Hence $v\circ e\in\funs{J}{\carrier[M]}$ by
        \obv{funs_circ}.
    $v$ forms a basis of $M$ over $R$ indexed by $I$ iff
        $v$ spans $M$ over $R$ indexed by $I$ and
        $v$ is linearly independent in $M$ over $R$ indexed by $I$ by
        \obv{indexed_module_basis_iff}.
    $v$ spans $M$ over $R$ indexed by $I$ iff
        $v\circ e$ spans $M$ over $R$ indexed by $J$ by
        \obv{indexed_module_spanning_reindex}.
    $v$ is linearly independent in $M$ over $R$ indexed by $I$ iff
        $v\circ e$ is linearly independent in $M$ over $R$ indexed by $J$ by
        \obv{indexed_module_linear_independence_reindex}.
    $v\circ e$ forms a basis of $M$ over $R$ indexed by $J$ iff
        $v\circ e$ spans $M$ over $R$ indexed by $J$ and
        $v\circ e$ is linearly independent in $M$ over $R$ indexed by $J$ by
        \obv{indexed_module_basis_iff}.
    Follows by \auto.
\end{proof}

\subsection{Consequences for spanning and bases}

\begin{theorem}\label{indexed_module_spanning_iff_span_full}
    Let $R$ be a ring.
    Let $M$ be a left module over $R$.
    Let $v\in\funs{I}{\carrier[M]}$.
    Then $v$ spans $M$ over $R$ indexed by $I$ iff
        $\moduleSpan{R}{M}{\ran{v}}=\carrier[M]$.
\end{theorem}
\begin{proof}
    Let $S=\FreeModule{R}{I}$ and
        $f=\linearCombinationMap{R}{I}{M}{v}$.
    $f$ is a linear map from $S$ to $M$ over $R$ by
        \auto{linear_combination_map_linear}.
    $f\in\funs{\carrier[S]}{\carrier[M]}$ by
        \obv{linear_map_type}.
    Hence $\dom{f}=\carrier[S]$ by \obv{funs}.
    Hence $\ran{f}=\img{f}{\carrier[S]}$ by \auto{img_dom}.
    Hence $\ran{f}=\moduleImage{f}{S}$ by
        \obv{module_image}.
    $f$ is a module surjection from $S$ to $M$ over $R$ iff
        $\moduleImage{f}{S}=\carrier[M]$ by
        \obv{module_surjection_iff_image_full}.
    Hence $f$ is a module surjection from $S$ to $M$ over $R$ iff
        $\ran{f}=\carrier[M]$ by \auto.
    $\moduleSpan{R}{M}{\ran{v}}=\ran{f}$ by
        \auto{indexed_family_span_eq_linear_combination_range}.
    Follows by \auto{indexed_module_spanning}.
\end{proof}

\begin{theorem}\label{indexed_module_basis_coordinate_existence}
    Let $R$ be a ring.
    Let $M$ be a left module over $R$.
    Let $v\in\funs{I}{\carrier[M]}$.
    Suppose $v$ forms a basis of $M$ over $R$ indexed by $I$.
    Let $y\in\carrier[M]$.
    Then there exists $x\in\carrier[\FreeModule{R}{I}]$ such that
        $\apply{\linearCombinationMap{R}{I}{M}{v}}{x}=y$.
\end{theorem}
\begin{proof}
    Let $S=\FreeModule{R}{I}$ and
        $f=\linearCombinationMap{R}{I}{M}{v}$.
    $\linearCombinationMap{R}{I}{M}{v}$ is a module isomorphism from
        $\FreeModule{R}{I}$ to $M$ over $R$ by
        \obv{indexed_module_basis}.
    Hence $f$ is a module isomorphism from $S$ to $M$ over $R$.
    Hence $f\in\bijections{\carrier[S]}{\carrier[M]}$ by
        \obv{module_isomorphism_is_bijection}.
    Hence $f\in\surj{\carrier[S]}{\carrier[M]}$ by
        \obv{bijections}.
    Follows by \auto{surj}.
\end{proof}

\begin{theorem}\label{indexed_module_coordinate_uniqueness}
    Let $R$ be a ring.
    Let $M$ be a left module over $R$.
    Let $v\in\funs{I}{\carrier[M]}$.
    Suppose $v$ is linearly independent in $M$ over $R$ indexed by $I$.
    Let $x,y\in\carrier[\FreeModule{R}{I}]$.
    Suppose
        $\apply{\linearCombinationMap{R}{I}{M}{v}}{x}
            =\apply{\linearCombinationMap{R}{I}{M}{v}}{y}$.
    Then $x=y$.
\end{theorem}
\begin{proof}
    Let $S=\FreeModule{R}{I}$ and
        $f=\linearCombinationMap{R}{I}{M}{v}$.
    $f$ is a linear map from $S$ to $M$ over $R$ by
        \auto{linear_combination_map_linear}.
    Hence $f\in\funs{\carrier[S]}{\carrier[M]}$ by
        \obv{linear_map_type}.
    Hence $\dom{f}=\carrier[S]$ and $f$ is a function by
        \obv{funs,funs_is_function}.
    $\linearCombinationMap{R}{I}{M}{v}$ is injective by
        \obv{indexed_module_linear_independence}.
    Hence $f$ is injective.
    Follows by \auto{injective_function}.
\end{proof}

\begin{theorem}\label{linear_maps_equal_on_indexed_basis}
    Let $R$ be a ring.
    Let $M$ be a left module over $R$.
    Let $N$ be a left module over $R$.
    Let $v\in\funs{I}{\carrier[M]}$.
    Suppose $v$ forms a basis of $M$ over $R$ indexed by $I$.
    Suppose $f$ is a linear map from $M$ to $N$ over $R$ and
        $g$ is a linear map from $M$ to $N$ over $R$.
    Suppose for every $i\in I$ we have $f(v(i))=g(v(i))$.
    Then $f=g$.
\end{theorem}
\begin{proof}
    Let $A=\ran{v}$.
    $v$ spans $M$ over $R$ indexed by $I$ by
        \auto{indexed_module_basis_iff}.
    Hence $\moduleSpan{R}{M}{A}=\carrier[M]$ by
        \auto{indexed_module_spanning_iff_span_full}.
    $A\subseteq\carrier[M]$ by \auto{funs_ran}.
    Show that for every $a\in A$ we have $f(a)=g(a)$.
    \begin{subproof}
        Fix $a\in A$.
        $v$ is a function by \obv{funs_is_function}.
        Take $i\in\dom{v}$ such that $v(i)=a$ by
            \auto{fun_ran_iff}.
        Hence $i\in I$ by \auto{funs}.
        Follows by \auto.
    \end{subproof}
    Follows by \auto{linear_map_equality_from_spanning_set}.
\end{proof}

\subsection{Phrase fixtures}

The following elementary statements keep the intended singular, plural, and
negated readings of the three global notions in the maintained library.

\begin{proposition}\label{indexed_module_spanning_phrase_fixture}
    Let $v$ be a set that spans $M$ over $R$ indexed by $I$.
    Then $v$ spans $M$ over $R$ indexed by $I$ and it is not the case that
        $v$ does not span $M$ over $R$ indexed by $I$.
\end{proposition}
\begin{proof}
\end{proof}

\begin{proposition}\label{indexed_module_linear_independence_phrase_fixture}
    Let $v$ be a set that is linearly independent in $M$ over $R$
        indexed by $I$.
    Then $v$ is linearly independent in $M$ over $R$ indexed by $I$ and
        it is not the case that
            $v$ is not linearly independent in $M$ over $R$ indexed by $I$.
\end{proposition}
\begin{proof}
\end{proof}

\begin{proposition}\label{indexed_module_basis_phrase_fixture}
    Let $v$ be a set that forms a basis of $M$ over $R$ indexed by $I$.
    Then $v$ forms a basis of $M$ over $R$ indexed by $I$ and
        it is not the case that
            $v$ does not form a basis of $M$ over $R$ indexed by $I$.
\end{proposition}
\begin{proof}
\end{proof}

\begin{proposition}\label{indexed_module_spanning_plural_phrase_fixture}
    Suppose $v$, $w$ span $M$ over $R$ indexed by $I$.
    Then $v$ spans $M$ over $R$ indexed by $I$ and
        $w$ spans $M$ over $R$ indexed by $I$.
\end{proposition}
\begin{proof}
\end{proof}

\begin{proposition}\label{indexed_module_basis_plural_phrase_fixture}
    Suppose $v$, $w$ form a basis of $M$ over $R$ indexed by $I$.
    Then $v$ forms a basis of $M$ over $R$ indexed by $I$ and
        $w$ forms a basis of $M$ over $R$ indexed by $I$.
\end{proposition}
\begin{proof}
\end{proof}
